% CS771 Progress Report -- uses the CVPR author kit (https://github.com/cvpr-org/author-kit/releases)
% Put this file next to cvpr.sty (and the other kit files) from the release zip, or upload into the kit on Overleaf.
% Prof requires 11pt, double column, final/camera-ready mode.
\documentclass[11pt,twocolumn,letterpaper]{article}
\usepackage{cvpr}              % no [review] option => final / camera-ready version
\usepackage{graphicx}
\usepackage{amsmath}
\usepackage{booktabs}
\usepackage{xcolor}
\usepackage[pagebackref,breaklinks,colorlinks,allcolors=cvprblue]{hyperref}

% Highlight anything that still needs real content. Delete before submitting.
\newcommand{\todo}[1]{\textcolor{red}{[TODO: #1]}}

\def\paperID{****}
\def\confName{CS771}
\def\confYear{2026}

\title{Occlusion-Robust Spatial-Temporal Graph Convolutional Networks\\for Smart Environment Action Recognition\\{\large Progress Report}}

\author{Karthik Reddy Jannupalli \and Sanjana Kallingal \and Penchala Siva Prasanth Kasa}

\begin{document}
\maketitle

%--------------------------------------------------------------------
\section{Current Status and Scope Change}
We are building an occlusion-robust skeleton action recognizer on ST-GCN~\cite{yan2018stgcn}
for smart-home settings, using a domestic-action subset of NTU RGB+D~\cite{shahroudy2016ntu}.
\textbf{Change from the proposal:} following the feedback, we have replaced the fixed
lower-body cluster dropout with a \emph{geometry-driven} occlusion model that covers both
object occlusion and self-occlusion. We have also redesigned the evaluation so that
test-time occlusion differs from training-time augmentation and includes naturally
occluded real-home data. The edge-efficiency study is now secondary.
\todo{confirm/adjust: which of these changes the team is actually adopting}

%--------------------------------------------------------------------
\section{Work Completed}
\textbf{Data pipeline.} \todo{e.g., selected N domestic classes (list a few), cross-subject/cross-view splits, 25-joint Kinect graph; status of MediaPipe~\cite{lugaresi2019mediapipe}/YOLOv8-Pose~\cite{jocher2023yolov8} 2D keypoint extraction and joint mapping.}

\textbf{Models.} \todo{which of Bi-LSTM, 1D-CNN, ST-GCN are implemented and trained; framework and hardware; training epochs.}

\textbf{Occlusion simulation.} Beyond whole-cluster dropout, we model
(i) \emph{object occlusion} by placing virtual furniture-like occluders (boxes of varying
height and position) in 3D and masking joints whose camera ray intersects the occluder;
(ii) \emph{self-occlusion} by using NTU's three camera views and joint depth relative to
the torso to mark hidden joints; (iii) \emph{temporal occlusion}, where occluders cover
contiguous frame blocks; and (iv) \emph{estimator-like corruption}, where occluded joints
receive jitter or held last-seen positions instead of clean zeros.
\todo{mark which of (i)--(iv) are implemented vs.\ planned.}

\textbf{Preliminary results.} \todo{Fill the table with real numbers; drop rows that are not run yet.}

\begin{table}[h]
\centering
\small
\begin{tabular}{@{}lccc@{}}
\toprule
Model & Clean & Lower-body occ. & Unseen occ. \\
\midrule
Bi-LSTM              & -- & -- & -- \\
1D-CNN               & -- & -- & -- \\
ST-GCN (no aug.)     & -- & -- & -- \\
ST-GCN + rand.\ joint drop & -- & -- & -- \\
ST-GCN + ours        & -- & -- & -- \\
\bottomrule
\end{tabular}
\caption{Top-1 accuracy (\%) on the NTU domestic subset. \todo{define each test condition precisely.}}
\label{tab:prelim}
\end{table}

\textbf{Difficulties.} \todo{real issues, e.g., NTU download size/preprocessing time, 2D--3D joint-topology mismatch between MediaPipe (33) / COCO (17) / Kinect (25), GPU limits, Kinect joints under occlusion are noisy rather than missing.}

%--------------------------------------------------------------------
\section{Next Steps}
\begin{enumerate}\setlength{\itemsep}{0pt}
\item \textbf{Finish the occlusion model} \todo{remaining parts} and confidence-aware input (joint confidence as an extra channel).
\item \textbf{Held-out occlusion evaluation:} train with our augmentation, then test on occlusion types not used in training (arms, torso, random regions, frame blocks), following the part- and frame-level occlusion protocol of RA-GCN~\cite{song2020ragcn}.
\item \textbf{Real occlusion:} evaluate on naturally occluded data, namely the Toyota Smarthome dataset~\cite{das2019toyota} (real apartments, furniture occlusion) and NTU samples with low Kinect tracking confidence.
\item \textbf{Ablation:} no augmentation vs.\ random-joint dropout vs.\ cluster dropout vs.\ geometric occlusion.
\item \textbf{Efficiency:} latency and parameter count on edge hardware (if time permits).
\end{enumerate}
\todo{add owner per item and rough dates up to the final deadline.}

%--------------------------------------------------------------------
\section{Response to Proposal Feedback}
\textbf{Comment:} random masking may not simulate realistic self- and object occlusion.
\textbf{Response:} our proposal already masked anatomical clusters rather than random joints,
but we agree that a fixed cluster dropped at random does not reflect real occlusion. We now derive
masks from scene geometry (virtual occluders and camera viewpoint), add self-occlusion from
multi-view depth, make occlusion temporally contiguous, and corrupt occluded joints the way
pose estimators do instead of zeroing them (Sec.~2).

\textbf{Comment:} the evaluation does not consider realistic occlusion, so generalization is unclear.
\textbf{Response:} previously, the test-time occlusion matched the training augmentation.
We now (a) test on occlusion types unseen during training, (b) evaluate on naturally occluded
real-home data (Toyota Smarthome) and low-confidence NTU samples, and (c) add an ablation
isolating the effect of the occlusion model (Sec.~3).

{\small
\bibliographystyle{ieeenat_fullname}
\begin{thebibliography}{9}
\bibitem{yan2018stgcn} S.~Yan, Y.~Xiong, and D.~Lin. Spatial temporal graph convolutional networks for skeleton-based action recognition. In \emph{AAAI}, 2018.
\bibitem{shahroudy2016ntu} A.~Shahroudy, J.~Liu, T.-T.~Ng, and G.~Wang. NTU RGB+D: A large scale dataset for 3D human activity analysis. In \emph{CVPR}, 2016.
\bibitem{lugaresi2019mediapipe} C.~Lugaresi et al. MediaPipe: A framework for building perception pipelines. \emph{arXiv:1906.08172}, 2019.
\bibitem{jocher2023yolov8} G.~Jocher, A.~Chaurasia, and J.~Qiu. Ultralytics YOLOv8. \url{https://github.com/ultralytics/ultralytics}, 2023.
\bibitem{song2020ragcn} Y.-F.~Song, Z.~Zhang, C.~Shan, and L.~Wang. Richly activated graph convolutional network for robust skeleton-based action recognition. \emph{IEEE TCSVT}, 2020.
\bibitem{das2019toyota} S.~Das, R.~Dai, M.~Koperski, L.~Minciullo, L.~Garattoni, F.~Bremond, and G.~Francesca. Toyota Smarthome: Real-world activities of daily living. In \emph{ICCV}, 2019.
\end{thebibliography}
}

\end{document}
